\documentclass[a4paper,11pt]{article}
\usepackage[utf8]{inputenc}
\usepackage[T1]{fontenc}
\usepackage{geometry}
\usepackage{enumitem}
\usepackage{titlesec}
\usepackage[table]{xcolor}
\usepackage{fancyhdr}
\usepackage{tabularx}
\usepackage{graphicx}
\usepackage{lastpage}
\usepackage{hyperref}
\usepackage{float}

% Page Setup
\geometry{top=1in, bottom=1in, left=0.8in, right=0.8in}
\pagestyle{fancy}
\fancyhf{}
\lhead{\small Project 83: Agentic AI Compiler}
\rhead{\small SRS Document}
\cfoot{\thepage\ of \pageref{LastPage}}

% Formatting
\titleformat{\section}{\Large\bfseries\color{blue!70!black}}{\thesection.}{1em}{}[\titlerule]
\titleformat{\subsection}{\large\bfseries\color{blue!50!black}}{\thesubsection}{1em}{}

% Custom commands
\newcommand{\priority}[1]{\colorbox{red!20}{\textbf{#1}}}
\newcommand{\prioritymed}[1]{\colorbox{yellow!20}{\textbf{#1}}}
\newcommand{\prioritylow}[1]{\colorbox{green!20}{\textbf{#1}}}

\begin{document}

\begin{center}
    {\Huge \textbf{Software Requirements}} \\
    {\Huge \textbf{Specification (SRS)}} \\
    \vspace{5mm}
    {\LARGE \textbf{Project 83}} \\
    {\large Conversational Agentic AI Compiler Assistant} \\
    {\large for Mini-Python Language} \\
    \vspace{5mm}
    \textbf{Version:} 1.0 \\
    \textbf{Date:} January 2026 \\
    \textbf{Status:} APPROVED
\end{center}

\vspace{5mm}

\tableofcontents
\clearpage

\section{Introduction}

\subsection{1.1 Purpose}

This Software Requirements Specification (SRS) document provides a complete description of all functionalities, constraints, and interfaces for the \textbf{Agentic AI Compiler Assistant} system. It is intended for:

\begin{itemize}
    \item \textbf{Developers:} To guide implementation of system components
    \item \textbf{Testers:} To create comprehensive test plans
    \item \textbf{Stakeholders:} To understand system capabilities and limitations
    \item \textbf{Academic Reviewers:} To evaluate project scope and complexity
\end{itemize}

\subsection{1.2 Scope}

\textbf{Product Name:} Agentic AI Compiler Assistant (Project 83)

\textbf{Product Description:} An intelligent compiler for the Mini-Python programming language that combines traditional compilation phases (lexical analysis, syntax parsing, semantic analysis) with AI-powered assistance to provide contextual error explanations, correction suggestions, and security auditing.

\textbf{Benefits:}
\begin{itemize}
    \item Reduces debugging time through intelligent error explanations
    \item Educates users on compiler theory and language grammar
    \item Identifies security vulnerabilities proactively
    \item Provides interactive learning environment for compiler concepts
\end{itemize}

\textbf{Goals:}
\begin{enumerate}
    \item Create a working compiler front-end for Mini-Python
    \item Integrate Google Gemini AI for intelligent assistance
    \item Achieve >95\% test coverage on all components
    \item Provide both CLI and Web UI interfaces
    \item Complete implementation within 16 weeks
\end{enumerate}

\subsection{1.3 Definitions, Acronyms, and Abbreviations}

\begin{table}[H]
\centering
\begin{tabularx}{\textwidth}{|l|X|}
\hline
\rowcolor{blue!10}
\textbf{Term} & \textbf{Definition} \\
\hline
AI & Artificial Intelligence \\
\hline
API & Application Programming Interface \\
\hline
AST & Abstract Syntax Tree - hierarchical representation of source code \\
\hline
CLI & Command Line Interface \\
\hline
EBNF & Extended Backus-Naur Form - notation for grammar specification \\
\hline
EOF & End of File \\
\hline
Gemini & Google's Large Language Model API \\
\hline
Lexer & Lexical Analyzer - converts source code to tokens \\
\hline
Mini-Python & Subset of Python language used in this project \\
\hline
Parser & Syntax Analyzer - converts tokens to AST \\
\hline
SRS & Software Requirements Specification \\
\hline
Token & Atomic unit of source code (e.g., keyword, operator) \\
\hline
UI & User Interface \\
\hline
\end{tabularx}
\caption{Terminology}
\end{table}

\subsection{1.4 References}

\begin{enumerate}
    \item Python Language Reference: \url{https://docs.python.org/3/reference/}
    \item Google Gemini API Documentation: \url{https://ai.google.dev/docs}
    \item Compiler Design Principles (Aho, Sethi, Ullman)
    \item Mini-Python EBNF Grammar (Project Documentation)
\end{enumerate}

\subsection{1.5 Overview}

This SRS is organized into the following sections:
\begin{itemize}
    \item \textbf{Section 2:} Overall product description and context
    \item \textbf{Section 3:} Detailed functional requirements
    \item \textbf{Section 4:} External interface requirements
    \item \textbf{Section 5:} Non-functional requirements
    \item \textbf{Section 6:} System constraints and assumptions
\end{itemize}

\clearpage

\section{Overall Description}

\subsection{2.1 Product Perspective}

The Agentic AI Compiler Assistant is a \textbf{standalone system} with external dependencies on:
\begin{itemize}
    \item \textbf{Google Gemini API:} Cloud-based AI service
    \item \textbf{Python Runtime:} Execution environment
    \item \textbf{Operating System:} Windows/Linux/macOS
\end{itemize}

\textbf{System Context Diagram:}
\begin{center}
\begin{tikzpicture}[node distance=2cm, auto]
    \tikzstyle{block} = [rectangle, draw, fill=blue!20, text width=5em, text centered, minimum height=3em]
    \tikzstyle{cloud} = [draw, ellipse, fill=green!20, minimum height=2em]
    
    \node [block] (system) {Agentic AI Compiler};
    \node [cloud, left of=system, node distance=4cm] (user) {User};
    \node [cloud, right of=system, node distance=4cm] (gemini) {Gemini API};
    \node [cloud, below of=system] (files) {Source Files};
    
    \draw [->] (user) -- node {Code Input} (system);
    \draw [->] (system) -- node {Results} (user);
    \draw [->] (system) -- node {AI Query} (gemini);
    \draw [->] (gemini) -- node {Response} (system);
    \draw [->] (files) -- node {Read} (system);
\end{tikzpicture}
\end{center}

\subsection{2.2 Product Functions}

\textbf{Core Functions:}
\begin{enumerate}
    \item \textbf{Lexical Analysis:} Convert source code to token stream
    \item \textbf{Syntax Parsing:} Build Abstract Syntax Tree from tokens
    \item \textbf{Semantic Analysis:} Validate variable scopes and types
    \item \textbf{Error Detection:} Identify compilation errors at all phases
    \item \textbf{AI Assistance:} Generate intelligent error explanations
    \item \textbf{Security Auditing:} Flag insecure coding patterns
    \item \textbf{Interactive Debugging:} Multi-turn conversation for fixes
    \item \textbf{Code Visualization:} Display AST structure graphically
\end{enumerate}

\subsection{2.3 User Classes and Characteristics}

\begin{table}[H]
\centering
\begin{tabularx}{\textwidth}{|l|X|X|}
\hline
\rowcolor{blue!10}
\textbf{User Class} & \textbf{Characteristics} & \textbf{Technical Expertise} \\
\hline
Students & Learning compiler design & Beginner to Intermediate \\
\hline
Educators & Teaching compiler theory & Advanced \\
\hline
Developers & Prototyping Mini-Python code & Intermediate \\
\hline
Researchers & Studying AI in compilers & Advanced \\
\hline
\end{tabularx}
\caption{User Classes}
\end{table}

\subsection{2.4 Operating Environment}

\textbf{Hardware:}
\begin{itemize}
    \item Processor: x86-64 or ARM64
    \item RAM: Minimum 4GB, Recommended 8GB
    \item Storage: 500MB free space
    \item Network: Internet connection for AI features
\end{itemize}

\textbf{Software:}
\begin{itemize}
    \item OS: Windows 10+, macOS 11+, Linux (Ubuntu 20.04+)
    \item Python: Version 3.10 or higher
    \item Browser: Chrome/Firefox (for Web UI)
\end{itemize}

\subsection{2.5 Design and Implementation Constraints}

\begin{itemize}
    \item \textbf{Language Constraint:} Must be implemented in Python 3.10+
    \item \textbf{No External Parsers:} Cannot use ANTLR, PLY, or similar tools
    \item \textbf{AI Dependency:} Requires active Gemini API key
    \item \textbf{Grammar Limitation:} Restricted to Mini-Python subset
    \item \textbf{API Rate Limits:} Subject to Gemini free tier quotas
\end{itemize}

\subsection{2.6 Assumptions and Dependencies}

\textbf{Assumptions:}
\begin{itemize}
    \item Users have basic understanding of Python syntax
    \item Users can obtain Google Gemini API key
    \item Internet connectivity is available for AI features
    \item Source code files are UTF-8 encoded
\end{itemize}

\textbf{Dependencies:}
\begin{itemize}
    \item \texttt{google-generativeai} Python package
    \item \texttt{python-dotenv} for configuration
    \item \texttt{pytest} for testing
    \item \texttt{streamlit} for Web UI (optional)
\end{itemize}

\clearpage

\section{Functional Requirements}

\subsection{3.1 Lexical Analysis Module}

\textbf{FR1.1: Token Generation} \priority{HIGH}

\textbf{Description:} The system shall convert Mini-Python source code into a sequence of tokens.

\textbf{Input:} String containing Mini-Python source code

\textbf{Output:} List of Token objects, each containing:
\begin{itemize}
    \item Token type (e.g., KEYWORD, IDENTIFIER, OPERATOR)
    \item Token value (e.g., "def", "x", "+")
    \item Line number
    \item Column number
\end{itemize}

\textbf{Processing Rules:}
\begin{enumerate}
    \item Recognize all keywords: def, if, elif, else, while, for, return, in, and, or, not, True, False
    \item Identify operators: +, -, *, /, \%, ==, !=, <, >, <=, >=
    \item Parse literals: integers, strings (single/double quotes)
    \item Extract identifiers: alphanumeric + underscore, starting with letter
    \item Handle delimiters: parentheses, colon, comma
\end{enumerate}

\vspace{3mm}
\textbf{FR1.2: Indentation Handling} \priority{HIGH}

\textbf{Description:} The system shall manage Python-style significant whitespace by emitting INDENT and DEDENT tokens.

\textbf{Processing Rules:}
\begin{enumerate}
    \item Maintain indentation stack initialized with [0]
    \item Emit INDENT when indentation increases
    \item Emit DEDENT(s) when indentation decreases
    \item Raise IndentationError for misaligned dedents
    \item Ignore blank lines and comments for indentation calculation
\end{enumerate}

\vspace{3mm}
\textbf{FR1.3: Lexical Error Detection} \priority{HIGH}

\textbf{Description:} The system shall detect and report lexical errors.

\textbf{Error Types:}
\begin{itemize}
    \item Invalid characters (e.g., \texttt{@}, \texttt{\$})
    \item Unterminated string literals
    \item Mixed tabs and spaces in indentation
    \item Invalid numeric literals
\end{itemize}

\textbf{Output:} Error message with line, column, and description

\subsection{3.2 Syntax Parsing Module}

\textbf{FR2.1: AST Construction} \priority{HIGH}

\textbf{Description:} The system shall build an Abstract Syntax Tree from the token stream according to Mini-Python grammar.

\textbf{Input:} Token list from Lexer

\textbf{Output:} AST root node representing program structure

\textbf{AST Node Types:}
\begin{itemize}
    \item ProgramNode
    \item FunctionDefNode
    \item IfStatementNode
    \item WhileStatementNode
    \item AssignmentNode
    \item BinaryOpNode
    \item VariableNode
    \item LiteralNode
\end{itemize}

\vspace{3mm}
\textbf{FR2.2: Syntax Error Detection} \priority{HIGH}

\textbf{Description:} The system shall detect violations of Mini-Python grammar.

\textbf{Error Types:}
\begin{itemize}
    \item Missing colons after function/if/while headers
    \item Mismatched parentheses
    \item Invalid expression syntax
    \item Unexpected tokens
\end{itemize}

\vspace{3mm}
\textbf{FR2.3: Operator Precedence} \priority{MEDIUM}

\textbf{Description:} The system shall enforce correct operator precedence and associativity.

\textbf{Precedence Order (High to Low):}
\begin{enumerate}
    \item Parentheses: \texttt{( )}
    \item Multiplication, Division, Modulo: \texttt{*}, \texttt{/}, \texttt{\%}
    \item Addition, Subtraction: \texttt{+}, \texttt{-}
    \item Comparison: \texttt{<}, \texttt{>}, \texttt{<=}, \texttt{>=}, \texttt{==}, \texttt{!=}
    \item Logical: \texttt{and}, \texttt{or}, \texttt{not}
\end{enumerate}

\subsection{3.3 Semantic Analysis Module}

\textbf{FR3.1: Symbol Table Management} \priority{HIGH}

\textbf{Description:} The system shall maintain a symbol table to track variable and function declarations.

\textbf{Functionality:}
\begin{itemize}
    \item Record variable declarations with scope
    \item Track function definitions and parameters
    \item Detect undeclared variable usage
    \item Identify variable redeclaration errors
\end{itemize}

\vspace{3mm}
\textbf{FR3.2: Scope Analysis} \priority{HIGH}

\textbf{Description:} The system shall enforce scope rules for variables and functions.

\textbf{Scope Rules:}
\begin{itemize}
    \item Global scope: Module-level definitions
    \item Local scope: Function parameters and local variables
    \item Nested scope: Variables from outer functions
\end{itemize}

\vspace{3mm}
\textbf{FR3.3: Type Checking} \priority{MEDIUM}

\textbf{Description:} The system shall perform basic type validation.

\textbf{Type Checks:}
\begin{itemize}
    \item Integer operations require integer operands
    \item Comparison operators require compatible types
    \item Function calls match declared parameters
\end{itemize}

\subsection{3.4 AI Assistance Module}

\textbf{FR4.1: Error Context Generation} \priority{HIGH}

\textbf{Description:} Upon compilation error, the system shall prepare detailed context for AI analysis.

\textbf{Context Contents:}
\begin{itemize}
    \item Error type and message
    \item Problematic code snippet (±5 lines)
    \item Token stream around error location
    \item Current AST structure (if available)
    \item Relevant grammar rules from EBNF
\end{itemize}

\vspace{3mm}
\textbf{FR4.2: AI Query Execution} \priority{HIGH}

\textbf{Description:} The system shall send error context to Gemini API and retrieve response.

\textbf{Processing:}
\begin{enumerate}
    \item Format error context as natural language prompt
    \item Inject EBNF grammar as system instruction
    \item Call Gemini API with configured model
    \item Parse AI response for explanation and suggestions
\end{enumerate}

\textbf{API Requirements:}
\begin{itemize}
    \item Use system instructions for "Compiler Assistant" persona
    \item Configure temperature=0.3 for consistent responses
    \item Handle API errors gracefully (timeout, rate limit)
\end{itemize}

\vspace{3mm}
\textbf{FR4.3: Response Formatting} \priority{MEDIUM}

\textbf{Description:} The system shall present AI responses in user-friendly format.

\textbf{Response Structure:}
\begin{itemize}
    \item \textbf{Error Explanation:} What went wrong and why
    \item \textbf{Grammar Reference:} Relevant EBNF rule
    \item \textbf{Suggested Fix:} Corrected code snippet
    \item \textbf{Additional Tips:} Best practices (optional)
\end{itemize}

\vspace{3mm}
\textbf{FR4.4: Conversational Loop} \priority{MEDIUM}

\textbf{Description:} The system shall support multi-turn dialogue for iterative debugging.

\textbf{Functionality:}
\begin{itemize}
    \item Maintain conversation history
    \item Allow user to ask follow-up questions
    \item Reference previous error context in new responses
    \item Option to reset conversation
\end{itemize}

\subsection{3.5 Security Auditing Module}

\textbf{FR5.1: Pattern Detection} \priority{HIGH}

\textbf{Description:} The system shall identify insecure coding patterns in Mini-Python code.

\textbf{Dangerous Patterns:}
\begin{itemize}
    \item Use of \texttt{eval()} function
    \item Use of \texttt{exec()} function
    \item String concatenation in loops (performance issue)
    \item Unused variables (code smell)
\end{itemize}

\vspace{3mm}
\textbf{FR5.2: Security Warnings} \priority{MEDIUM}

\textbf{Description:} The system shall generate warnings for detected security risks.

\textbf{Warning Format:}
\begin{itemize}
    \item Severity level (HIGH, MEDIUM, LOW)
    \item Description of risk
    \item Secure alternative suggestion
    \item Reference to security best practices
\end{itemize}

\subsection{3.6 User Interface Module}

\textbf{FR6.1: Command Line Interface} \priority{HIGH}

\textbf{Description:} The system shall provide a CLI for code compilation and interactive debugging.

\textbf{Commands:}
\begin{itemize}
    \item \texttt{compile <filename>}: Compile source file
    \item \texttt{debug}: Enter interactive debugging mode
    \item \texttt{help}: Display usage instructions
    \item \texttt{exit}: Quit application
\end{itemize}

\vspace{3mm}
\textbf{FR6.2: Web User Interface} \prioritymed{MEDIUM}

\textbf{Description:} The system shall provide a Streamlit-based web interface.

\textbf{UI Components:}
\begin{itemize}
    \item Left panel: Code editor with syntax highlighting
    \item Right panel: AI chatbot interface
    \item Bottom panel: AST visualization
    \item Toolbar: Compile, Clear, Download buttons
\end{itemize}

\vspace{3mm}
\textbf{FR6.3: AST Visualization} \prioritylow{LOW}

\textbf{Description:} The system shall display Abstract Syntax Tree graphically.

\textbf{Visualization Features:}
\begin{itemize}
    \item Tree structure with expandable nodes
    \item Node details on hover/click
    \item Export as PNG/SVG
\end{itemize}

\clearpage

\section{External Interface Requirements}

\subsection{4.1 User Interfaces}

\textbf{UI1: Command Line Interface}
\begin{itemize}
    \item Standard terminal/console application
    \item Text-based input and output
    \item Color-coded error messages (red) and success (green)
    \item Supports ANSI escape codes for formatting
\end{itemize}

\textbf{UI2: Web User Interface}
\begin{itemize}
    \item Browser-based interface using Streamlit
    \item Responsive design (desktop and tablet)
    \item Minimum screen resolution: 1280x720
    \item Dark mode support
\end{itemize}

\subsection{4.2 Hardware Interfaces}

No direct hardware interfaces required. System uses standard OS I/O for:
\begin{itemize}
    \item Keyboard input
    \item Display output
    \item File system access
\end{itemize}

\subsection{4.3 Software Interfaces}

\textbf{SI1: Google Gemini API}
\begin{itemize}
    \item Interface: RESTful HTTP API
    \item Protocol: HTTPS
    \item Authentication: API Key in request headers
    \item Data Format: JSON
    \item Models: gemini-1.5-pro, gemini-1.5-flash, or later
\end{itemize}

\textbf{SI2: Python Standard Library}
\begin{itemize}
    \item \texttt{re} - Regular expression operations
    \item \texttt{os} - Operating system interface
    \item \texttt{sys} - System-specific parameters
    \item \texttt{enum} - Enumeration support
    \item \texttt{typing} - Type hints
\end{itemize}

\textbf{SI3: Third-Party Libraries}
\begin{itemize}
    \item \texttt{google-generativeai} - Gemini SDK
    \item \texttt{python-dotenv} - Environment variables
    \item \texttt{pytest} - Testing framework
    \item \texttt{streamlit} - Web UI framework
\end{itemize}

\subsection{4.4 Communication Interfaces}

\textbf{Network Communication:}
\begin{itemize}
    \item Protocol: HTTPS (TLS 1.2+)
    \item Port: 443 (standard HTTPS)
    \item Endpoint: \texttt{generativelanguage.googleapis.com}
    \item Data: JSON payloads
\end{itemize}

\textbf{File I/O:}
\begin{itemize}
    \item Input: \texttt{.py} files (UTF-8 encoding)
    \item Output: Console logs, JSON reports
    \item Configuration: \texttt{.env} files
\end{itemize}

\clearpage

\section{Non-Functional Requirements}

\subsection{5.1 Performance Requirements}

\textbf{NFR1.1: Compilation Speed} \priority{HIGH}

\begin{itemize}
    \item Lexer: Process at least 1,000 lines/second
    \item Parser: Build AST in <200ms for typical programs (<200 lines)
    \item Semantic Analysis: Complete in <100ms
    \item Total local processing: <500ms for 200-line program
\end{itemize}

\textbf{NFR1.2: AI Response Time} \prioritymed{MEDIUM}

\begin{itemize}
    \item Gemini API call: <3 seconds (90th percentile)
    \item Total error-to-suggestion latency: <5 seconds
    \item Timeout handling: 10 seconds maximum wait
\end{itemize}

\textbf{NFR1.3: Memory Usage} \prioritylow{LOW}

\begin{itemize}
    \item Token storage: <100 bytes per token
    \item AST nodes: <200 bytes per node
    \item Total memory for 500-line program: <5 MB
\end{itemize}

\subsection{5.2 Safety Requirements}

\textbf{NFR2.1: API Key Security} \priority{HIGH}

\begin{itemize}
    \item API keys stored in \texttt{.env} file (not version controlled)
    \item Keys loaded via \texttt{python-dotenv} at runtime
    \item No hardcoding of credentials in source code
    \item Clear error message if key is missing/invalid
\end{itemize}

\textbf{NFR2.2: Error Recovery} \prioritymed{MEDIUM}

\begin{itemize}
    \item System shall not crash on malformed input
    \item Graceful handling of API failures (timeout, network error)
    \item Comprehensive exception handling throughout
    \item User-friendly error messages (no stack traces to end-users)
\end{itemize}

\subsection{5.3 Security Requirements}

\textbf{NFR3.1: Input Validation} \priority{HIGH}

\begin{itemize}
    \item Sanitize file paths to prevent directory traversal
    \item Limit file size to 10 MB maximum
    \item Validate file encoding (UTF-8 only)
    \item Reject files with dangerous extensions
\end{itemize}

\textbf{NFR3.2: AI Prompt Injection Prevention} \priority{HIGH}

\begin{itemize}
    \item Escape user code before sending to AI
    \item Use system instructions to lock AI persona
    \item Validate AI responses for unexpected content
    \item Rate limiting on API calls (prevent abuse)
\end{itemize}

\subsection{5.4 Software Quality Attributes}

\textbf{NFR4.1: Reliability} \priority{HIGH}

\begin{itemize}
    \item \textbf{MTBF (Mean Time Between Failures):} >100 hours of operation
    \item \textbf{Test Coverage:} Minimum 95\% for core modules
    \item \textbf{Error Rate:} <1\% false positives in syntax detection
\end{itemize}

\textbf{NFR4.2: Maintainability} \prioritymed{MEDIUM}

\begin{itemize}
    \item \textbf{Code Documentation:} Docstrings for all public functions
    \item \textbf{Modularity:} Clear separation of Lexer, Parser, AI modules
    \item \textbf{Naming Conventions:} PEP 8 compliance
    \item \textbf{Cyclomatic Complexity:} <10 per function
\end{itemize}

\textbf{NFR4.3: Portability} \prioritymed{MEDIUM}

\begin{itemize}
    \item \textbf{Cross-Platform:} Windows, macOS, Linux support
    \item \textbf{Python Versions:} Compatible with 3.10, 3.11, 3.12
    \item \textbf{Virtual Environment:} Fully functional in \texttt{venv}
    \item \textbf{No OS-Specific Code:} Use \texttt{os.path} instead of hardcoded paths
\end{itemize}

\textbf{NFR4.4: Usability} \prioritymed{MEDIUM}

\begin{itemize}
    \item \textbf{Learning Curve:} New users productive within 30 minutes
    \item \textbf{Error Messages:} Clear, actionable guidance
    \item \textbf{Documentation:} README with setup and usage examples
    \item \textbf{Help System:} Built-in \texttt{--help} command
\end{itemize}

\textbf{NFR4.5: Scalability} \prioritylow{LOW}

\begin{itemize}
    \item Handle source files up to 1,000 lines
    \item Support up to 100 function definitions per file
    \item Maintain performance for nesting depth up to 10 levels
\end{itemize}

\textbf{NFR4.6: Testability} \priority{HIGH}

\begin{itemize}
    \item Unit tests for all modules using \texttt{pytest}
    \item Integration tests for end-to-end compilation
    \item Mock Gemini API for offline testing
    \item Automated CI/CD pipeline (future)
\end{itemize}

\clearpage

\section{System Constraints}

\subsection{6.1 Regulatory Policies}

\begin{itemize}
    \item \textbf{Academic Integrity:} Code must be original (no copy from existing compilers)
    \item \textbf{API Usage:} Comply with Google Gemini Terms of Service
    \item \textbf{Open Source:} May be released under MIT License
\end{itemize}

\subsection{6.2 Hardware Limitations}

\begin{itemize}
    \item System designed for single-machine deployment (no distributed computing)
    \item No GPU acceleration required
    \item Standard consumer hardware sufficient
\end{itemize}

\subsection{6.3 Interfaces to Other Applications}

\begin{itemize}
    \item No integration with IDEs (VSCode, PyCharm) in initial version
    \item Standalone application (not a library)
    \item Future: Plugin architecture for extensibility
\end{itemize}

\subsection{6.4 Parallel Operations}

\begin{itemize}
    \item Single-threaded compilation pipeline
    \item Sequential processing: Lexer → Parser → Semantic → AI
    \item Async API calls may be considered for performance
\end{itemize}

\subsection{6.5 Audit Functions}

\begin{itemize}
    \item Logging of all compilation attempts (optional)
    \item Anonymized telemetry for error patterns (future)
    \item User can disable all logging via configuration
\end{itemize}

\subsection{6.6 Control Functions}

\begin{itemize}
    \item Configuration via \texttt{.env} file
    \item CLI flags for verbose mode, debug mode
    \item No runtime modification of grammar rules
\end{itemize}

\subsection{6.7 Higher-Order Language Requirements}

\begin{itemize}
    \item All code written in Python 3.10+
    \item No C/C++ extensions
    \item Pure Python implementation for portability
\end{itemize}

\subsection{6.8 Signal Handshake Protocols}

\begin{itemize}
    \item HTTP/HTTPS for API communication
    \item No custom network protocols
\end{itemize}

\subsection{6.9 Reliability Requirements}

\begin{itemize}
    \item Graceful degradation if AI unavailable (show basic error messages)
    \item Retry logic for transient API failures (max 3 retries)
    \item System remains functional without internet (compilation only, no AI)
\end{itemize}

\subsection{6.10 Criticality of the Application}

\begin{itemize}
    \item \textbf{Criticality Level:} LOW (educational/research tool)
    \item Not used in production environments
    \item Incorrect compilations do not pose safety risks
\end{itemize}

\subsection{6.11 Safety and Security Considerations}

\begin{itemize}
    \item Code execution not supported (analysis only)
    \item Sandbox not required (no runtime execution)
    \item API keys protected from accidental exposure
\end{itemize}

\clearpage

\section{Appendix A: Use Cases}

\subsection{Use Case 1: Student Debugging Syntax Error}

\textbf{Actor:} Computer Science Student

\textbf{Precondition:} Student has written Mini-Python code with syntax error

\textbf{Main Flow:}
\begin{enumerate}
    \item Student saves code to \texttt{factorial.py}
    \item Student runs \texttt{python main.py compile factorial.py}
    \item System detects missing colon after function definition
    \item System sends error context to Gemini API
    \item AI explains error and suggests fix
    \item Student corrects code and recompiles
    \item Compilation succeeds
\end{enumerate}

\textbf{Postcondition:} Student understands grammar rule and fixes error

\subsection{Use Case 2: Educator Demonstrating Compiler Phases}

\textbf{Actor:} University Professor

\textbf{Precondition:} Professor wants to show compilation pipeline in class

\textbf{Main Flow:}
\begin{enumerate}
    \item Professor opens Web UI in browser
    \item Types sample code with intentional errors
    \item Clicks "Compile" button
    \item AST visualization shows partial tree before error
    \item Error explanation displayed in chatbot panel
    \item Professor demonstrates fix and recompiles
    \item Complete AST displayed
\end{enumerate}

\textbf{Postcondition:} Students visualize compilation stages

\subsection{Use Case 3: Security Researcher Finding Vulnerabilities}

\textbf{Actor:} Security Analyst

\textbf{Precondition:} Analyst testing security auditing feature

\textbf{Main Flow:}
\begin{enumerate}
    \item Analyst writes code using \texttt{eval()}
    \item Runs compiler with security mode enabled
    \item System flags \texttt{eval()} as HIGH risk
    \item AI suggests secure alternative (e.g., \texttt{ast.literal\_eval()})
    \item Analyst reviews suggestion and refactors code
\end{enumerate}

\textbf{Postcondition:} Code hardened against injection attacks

\section{Appendix B: Data Dictionary}

\begin{table}[H]
\centering
\begin{tabularx}{\textwidth}{|l|X|}
\hline
\rowcolor{blue!10}
\textbf{Data Element} & \textbf{Description} \\
\hline
Token & Object with type, value, line, column attributes \\
\hline
TokenType & Enumeration of all possible token categories \\
\hline
AST Node & Abstract base class for syntax tree elements \\
\hline
Symbol Table & Dictionary mapping identifiers to type/scope info \\
\hline
Error Context & Structured data sent to AI (code, tokens, grammar) \\
\hline
AI Response & Natural language explanation + suggested fix \\
\hline
\end{tabularx}
\caption{Key Data Elements}
\end{table}

\section{Appendix C: Change History}

\begin{table}[H]
\centering
\begin{tabularx}{\textwidth}{|l|l|X|}
\hline
\rowcolor{blue!10}
\textbf{Version} & \textbf{Date} & \textbf{Changes} \\
\hline
1.0 & Jan 2026 & Initial SRS document created \\
\hline
\end{tabularx}
\caption{Document Revision History}
\end{table}

\vspace{10mm}
\noindent\textbf{Approval Signatures:} \\
\vspace{5mm}

\noindent\textbf{Project Lead:} \rule{5cm}{0.4pt} \hspace{1cm} \textbf{Date:} \rule{3cm}{0.4pt} \\
\vspace{5mm}

\noindent\textbf{Technical Reviewer:} \rule{5cm}{0.4pt} \hspace{1cm} \textbf{Date:} \rule{3cm}{0.4pt} \\

\end{document}
